\subsection{Five objects with cyclic support at nineteen arrows}

\begin{theorem}[Proved: cyclic five-object support through nineteen]
\label{CP2:A-strong-five-nineteen}
A finite category on five objects with cyclic support and at most
nineteen arrows cannot have classifying space of type $\CP^2$.
\end{theorem}
\begin{proof}
Assume a connected counterexample. Any actual arrow-reducing
homotopy equivalence, or any actual retract preserving projective-plane
type, contradicts Theorem~\ref{CP2:A-global-lower-nineteen}.
In particular no object is a retract of another. A strong component
on $t$ objects consequently has at least $t^2+t$ internal arrows.

A proper component on three objects costs at least twelve arrows.
If it is nonextreme, two interfaces cost at least six, their composite
costs one, and the two outside identities cost two, giving at least
twenty-one arrows. Thus it is extreme. Its two outside objects and
a connected interface cost at least six arrows, so its internal size
is at most thirteen. Its minimal group is trivial and its total row
and column counts are both three: a larger total would require at
least $3\cdot4+3=15$ internal arrows. There is one row and one column
at each object, so the actual source minimal-group collapse or its
sink dual strictly decreases the arrow count. Components of size at
least four are too expensive. If there are two strong pairs, first suppose there is a direct
interface between them. It costs at least four arrows, and a connected
attachment of the fifth object costs at least two more. Their two
internal costs and the fifth identity then give at least
$6+6+1+4+2=19$ arrows. Equality forces both pairs to be pure minimal
bands with their identities, and an extreme pair collapses strictly.
If there is no direct interface between the pairs, the fifth object
is their common source or their common sink: opposite directions
would create a pair-to-pair composite. The two interfaces cost at
least four arrows in total, so the two internal sizes have sum at
most fourteen. At least one pair has size at most seven. Its minimal
group is trivial and it has exactly one row and column at each
object; as an extreme component it collapses strictly. These cases
are excluded.

There remains one strong pair $S$, of internal size $q\ge6$, and
three singleton components. Let $D$ be the full three-object complement.
The elementary component count gives $q\le12$.

\emph{An extreme pair.} By opposition suppose $S$ is a source.
Let $\kappa$ be the number of connected components of $D$, and let
$h$ be the total interface cardinality. Then
\[
 |D|\ge6-\kappa,\qquad h\ge2\kappa,\qquad q+|D|+h\le19.
\]
If both block second cohomology groups vanish, use the two-block
product lemma. Suppose first that $H^2(BD;\mathbb Q)\ne0$.
A connected three-object positive-second-homology block costs at
least eight arrows. For two components the corresponding minimum
is seven, and for three components it is also seven. Thus $q\le9$.
If there is one minimal row per object, the source-component
collapse applies. Otherwise the budget leaves a trivial minimal
group, total rows three and total columns two. At $q=8$ the pair
is exactly its six-arrow band and two identities; its one-row
object satisfies the incoming one-sided deletion, with only the
identity as a unit. At $q=9$, a disconnected $D$ would cost at
least $7+4=11$ further arrows. Hence $D$ is connected, and its
minimum eight arrows forces $h=2$. Every external exit is therefore
singleton or empty, and the one-column fixed-exit collapse applies.
All these operations strictly reduce arrows.

If instead $H^2(BS;\mathbb Q)\ne0$, then $q=11$ or $12$,
the minimal group is trivial, and the total row and column counts
are both three. In particular each local row count is at most two.
If $\kappa\ge2$, the budget forces $q=11$, $\kappa=2$,
$|D|=4$ and $h=4$. The complement consists of trivial corners
and singleton hom-sets, and each nonempty interface is singleton.
The inclusion of $S$ factors through the join with either interface
target, hence is nullhomotopic; $H^2(BD;\mathbb Q)=0$ as well.
The two-block lemma applies.

For connected $D$, one has $|D|\ge5$. If $h=2$, both exits end
at one object $w$, giving the same nullhomotopy of the inclusion
of $S$. Also $|D|\le6$, so its second cohomology vanishes by
the homology cost bound. The only remaining possibility is
$q=11$, $|D|=5$, $h=3$. Now $D$ is a trivial-corner singleton
tree, and $S$ meets just one object $w$; otherwise four interface
arrows would be required. The two stalks into $w$ have sizes one
and two. The full subcategory $E=S\cup\{w\}$ has zero positive
products by the categorical binary-act lemma. No strict arrow
leaves $w$ inside $D$, since it would create a further interface.
Consequently the full nerve is $NE\cup ND$, meeting in the point
$w$. Degree-four restriction is injective, and both pieces have
zero second products. This concludes the extreme cases.

\emph{A nonextreme pair.} Choose outside objects $a,b$ with
$a<S<b$. There are four nonempty individual entrance and exit
hom-sets and the composite $a\to b$. Let $P$ be the total excess
over one in every nonempty external hom-set incident with $S$.
The same support count as in the ninth-excess proof gives
\[
 q+9+P\le19,\qquad 6\le q\le10,\qquad P\le10-q.
\]
Equality in this count means that all outside corners and outside
strict hom-sets are singleton, the contracted support is the triangle
$a,S,b$ with one additional leaf, and its only interfaces are
$a\to S$ and $S\to b$. The extra leaf cannot attach to $S$:
such an interface has two arrows instead of the one allowed by
equality. It therefore attaches at an outside trivial corner and
can be deleted by the singleton-comma lemma. Thus every equality
case $q+P=10$ is already excluded.

For $q=6$, the four local incoming/outgoing fixed-arrow conditions
of Lemma~\ref{CP2:A-high-one-sided} all hold internally; their
failure would require at least four external excesses. Hence
$P\le3$ gives a one-sided deletion, while $P=4$ is an equality
case. For $q=7$, the unique extra internal arrow can violate at
most two conditions. Thus $P\le1$ gives a deletion and $P=3$
is an equality case. Postpone $q=7,P=2$.

For $q=8$ or $9$, the minimal group is trivial. If both total
row and column counts are two, failure of both fixed-entrance
and fixed-exit collapses requires at least one external excess
on each side, so $P\ge2$. At $q=8$ this is an equality case;
at $q=9$ it is impossible. The other shapes are $(3,2)$ and
$(2,3)$. For $q=8$ the band is pure. In shape $(3,2)$, failure
of the fixed-exit collapse requires an exit excess. At the
one-row object, all internal entrances are fixed and the only
unit is the identity, so failure of its one-sided deletion
requires a further entrance excess. Again $P\ge2$ gives
equality. Shape $(2,3)$ is dual. At $q=9$, avoiding the
one-column or one-row block collapse requires $P\ge1$,
already an equality case. For $q=10$ one has $P=0$.
If the minimal group is trivial, one total is two and the
fixed-interface collapse applies. A nontrivial group forces
the pure $2\times2$ categorical band over $C_2$, plus two
global identities. A local minimal idempotent fixes all
internal incoming arrows and all singleton external entrances;
the only unit of its corner is its global identity. The
incoming one-sided deletion therefore applies here too.

It remains $q=7,P=2$. There is one minimal row and column
at each object. If all external entrances are singleton, the
fixed-entrance block collapse applies; if all external exits are
singleton, use the fixed-exit collapse. Avoiding both therefore
requires one external excess on each side. Thus each of the two
individual-stalk interfaces has one singleton and one doubled
stalk. No location for the unique extra internal arrow is needed. Only one arrow beyond the elementary
support budget remains. Let $c$ be the third outside object.
It cannot meet $S$: an additional interface together with
its forced composite to $a$ or $b$ costs two additional arrows
over a pendant edge, exceeding the available one. Thus $c$
meets the outside edge $a\to b$.

If it meets both $a$ and $b$, acyclicity and transitivity force
$a\to c\to b$, with trivial corners and singleton new hom-sets.
The unique singleton entrance from the trivial corner $a$
deletes $c$. If $c$ is pendant, transitivity forces its orientation
to be $a\to c$ or $c\to b$. A doubled pendant stalk uses the
single spare arrow, leaves both incident corners trivial, and
gives a nonzero first-homology circle. Otherwise its stalk is
singleton. If its neighbor is trivial, the singleton-comma
lemma deletes $c$, even if its own corner uses the spare arrow.
If its neighbor has order two, then $c$ is trivial. At $a$
there is a singleton entrance $a\to u$ for some $u\in S$,
and at $b$ there is a singleton exit $v\to b$ for some $v\in S$.
The actual singleton pendant-fold theorem therefore folds $c$
into $u$ or $v$, respectively, preserving projective-plane type while losing
arrows. This covers every placement of the spare arrow and
finishes the last case.
\end{proof}
